
Throughout this work, the experimental approach has progressively sharpened a single question: not simply whether a model gives the right answer, but where that answer comes from. The focus has been on what information the model actually has access to, and whether the different modalities are genuinely integrated during reasoning, rather than merely present.
The MAIA framework opens this line of inquiry. Its competence-oriented evaluation shows that high performance on individual items does not translate into robust multimodal understanding. When correctness is required to hold across minimally varied formulations and across both discriminative and generative tasks, the models become noticeably more fragile. The benchmark exposes a gap between producing a plausible-looking answer and maintaining a video representation that is grounded, consistent, and stable across different linguistic phrasings of the same underlying content. The difficulty models show with temporally distributed, relational, and implicit phenomena suggests that the core limitation goes beyond object recognition, extending to the construction of a systematic representation of events and how they relate to one another.
\textit{Lost in Translation} extends this line of inquiry from visual grounding to cross-modal interaction. Adding audio shows that the presence of an additional information channel does not by itself guarantee that a model exploits it. Audio can help when visual evidence is weak, but its effect shrinks once strong visual evidence is already available. Converting the same acoustic content into a transcription changes its effect on the model's behavior, so raw audio and its textual counterpart cannot be treated as interchangeable representations of the same information. This points to a basic property of current multimodal systems: how much a modality contributes depends on what other information is already available and on the representational form in which that additional evidence is supplied.
MAIA-AV takes this observation as its central object of study. The results show that multimodal reasoning is not simply the sum of independent unimodal capabilities. Audio, video, and transcription each supply evidence that partially overlaps but is not functionally identical. An additional modality can correct an error, leave a decision unchanged, or overturn a decision that was previously correct. Improved overall accuracy is therefore only one possible outcome of adding a modality; just as informative are the cases where a new modality changes a decision without any net gain, which show that integration involves competition between information sources as much as complementarity.
The unimodal conditions clarify this further. The strong contribution of audio on its own shows that speech, environmental sound, and their timing can carry information directly relevant to visually defined events. Transcription behaves differently, confirming that linguistic accessibility and perceptual depth are not the same property: once information is expressed in language, larger models draw on semantic associations more effectively, while increasing model size does little to resolve the difficulty of extracting information from visual input in the first place. This points to a bottleneck that sits before linguistic reasoning even begins, at the point where dynamic perceptual evidence must first be converted into the structured semantic representation the task requires.
The preliminary analysis supports the same interpretation. Models recover entities more reliably than they build complete representations of events and the relations between them. Evidence retrieved before the question is known correlates with downstream success but does not determine it: the question itself can direct the model's attention toward evidence it would otherwise miss. Correct task execution therefore depends on several stages, the initial availability of perceptual evidence, question-guided access to it, and its transformation into a representation suitable for comparison or inference, and a correct answer can still emerge even when one of these stages is only partially completed. This is why behavioral accuracy on its own cannot be taken as proof of genuine multimodal comprehension.
The no-input experiments confirm this from the opposite direction. Some decisions remain possible even with no perceptual evidence at all, showing that linguistic plausibility, world knowledge, and structural properties of the answer options can partly resolve the task by themselves. This does not make the perceptual component of the evaluation less important; it shows why it is necessary to separate what a model derives from the stimulus from what it already knows going in. A model should not be credited with genuine grounding simply because it reaches the correct answer, if that same answer is reachable through non-perceptual shortcuts.
The attention analysis adds a further layer to this picture. The models tested here do not distribute attention across audio, video, transcription, textual context, and their own generated history in the same way. The same modality can receive very different amounts of attention depending on the architecture and on the form in which the information is presented, and similar levels of attention can still lead to different behavioral outcomes. Attention therefore reflects access to information, but not whether that information is useful, discriminative, or causally responsible for the final decision. Evaluating multimodal integration means relating internal attention allocation to observed behavioral change, rather than treating either one as sufficient evidence on its own.
Together, these results trace a progression from \textit{multimodal availability} to \textit{multimodal integration}. MAIA showed that visual access and correct prediction alone do not guarantee robust grounding. \textit{Lost in Translation} showed that adding a modality does not guarantee that it contributes useful information. MAIA-AV shows that even when several modalities are each individually informative, using them together remains selective, asymmetric, and strongly shaped by architecture. The central limitation of current systems is accordingly not an inability to perceive each modality on its own, but a difficulty in building, selecting, and combining the information each one provides in a way that actually matches what the reasoning task requires.
Seen this way, genuine multimodal competence is not defined by how many input modalities a model can accept. It is defined by the ability to identify relevant evidence, keep complementary information from different representational spaces available, discard redundant or misleading signals, and combine what remains into a single coherent semantic representation. The results from MAIA-AV show that current models already have parts of this competence, even if putting them together is still far from complete. Multimodality, in this sense, is not a property of the input interface. It is a reasoning ability that has to be evaluated in its own right.
